\documentclass[../../../include/open-logic-section]{subfiles}

\begin{document}

\olfileid{sth}{ordinals}{basic}
\olsection{د ترتيبي عددونو بنسټيز خاصيتونه}

ومو ليدل چې لومړني څو ترتيبي عددونه طبيعي عددونه دي. د ترتيبي عددونو
د نظريې د جوړولو اصلي وجه دا ده چې د استقرا هغه اصل پراخ کړو چې په
طبيعي عددونو باندې عملي کېږي. دې مقصد ته به د ځينو ابتدايي نتيجو په
لړۍ ورسېږو.

\begin{lem}\ollabel{ordmemberord}
د ترتيبي عدد هر !!{element} ترتيبي عدد دے.
\end{lem}

\begin{proof}
فرض کړئ $\alpha$ ترتيبي عدد دے او $b \in \alpha$ دے. ځکه چې $\alpha$
متعدي دے، $b \subseteq \alpha$ دے. نو $\in$ د $b$ ښه ترتيب ورکوي، ځکه
چې $\in$ د $\alpha$ ښه ترتيب ورکوي.

د دې د ليدلو دپاره چې $b$ متعدي دے، فرض کړئ $x \in c \in b$ دي.
نو $c \in \alpha$ دے، ځکه چې $b
\subseteq \alpha$ دے. بيا، ځکه چې $\alpha$ متعدي دے، $c \subseteq
\alpha$ دے، نو $x \in \alpha$ دے. په دې توګه $x, c, b \in \alpha$ دي.
خو $\in$ د $\alpha$ ښه ترتيب ورکوي، نو د
\olref[sth][ordinals][wo]{wo:strictorder} له مخې $\in$ په
$\alpha$ باندې متعدي اړيکه ده. نو له $x
\in c \in b$ څخه $x \in b$ ترلاسه کوو. په تعميمولو سره $c \subseteq b$
\end{proof}

\begin{cor}\ollabel{ordissetofsmallerord}
$\alpha = \Setabs{\beta \in \alpha}{\beta \text{ ترتيبي عدد دے}}$، د
هر ترتيبي عدد~$\alpha$ دپاره.
\end{cor}

\begin{proof}
له \olref{ordmemberord} څخه سمدستي نتيجه اخلو.
\end{proof}

د لاندې دوو مهمو نتيجو،
\olref{ordinductionschema} او \olref{ordtrichotomy}، عمومي مطلب دا دے
چې ترتيبي عددونه پخپله د غړيتوب له مخې ښه مرتب دي:

\begin{thm}[د نامتناهيت هاخوا استقرا]\ollabel{ordinductionschema}
د هر فارمول $\phi(x)$ دپاره:
\[
	\text{که }\exists \alpha \phi(\alpha)\text{ وي، نو }\exists \alpha(\phi(\alpha)
	\land  (\forall \beta \in \alpha) \lnot \phi(\beta))
\]
دلته ښودل شوي کمّي کوونکي په ضمني توګه ترتيبي عددونو ته محدود دي.
\end{thm}

\begin{proof}
%\olref{ordinductionschema1}. 
فرض کړئ $\phi(\alpha)$ د کوم ترتيبي عدد $\alpha$ دپاره صدق کوي.
که $ (\forall \beta
\in \alpha) \lnot \phi(\beta)$ وي، نو مقصد ترلاسه شو. که داسې نۀ وي، نو ځکه
چې $\alpha$ ترتيبي عدد دے، د $\in$ له مخې يو تر ټولو کوچنے داسې !!{element}
لري چې $\phi$ پرې صدق کوي؛ او د \olref{ordmemberord} له مخې دا ترتيبي عدد دے.
%	\olref{ordinductionschema2}. Suppose there is some ordinal
%	$\gamma$ such that $\lnot\phi(\gamma)$. Then by
%	\olref{ordinductionschema1}, there is an $\in$-minimal ordinal
%	$\alpha$ for which $\lnot\phi(\alpha)$. So $(\forall \beta <
%	\alpha) \phi(\beta)$, rendering the antecedent of the conditional
%	false.
\end{proof}
\noindent
وګورئ چې \olref{ordinductionschema} په هماغه معنا د دې شېما په بڼه هم
بيانولے شو:
\[
\text{که }\forall \alpha((\forall \beta \in \alpha)\phi(\beta) \lif
\phi(\alpha))\text{ وي، نو }\forall \alpha\phi(\alpha)
\]
يوازې په \olref{ordinductionschema} کښې $\lnot\phi(\alpha)$ اخلو او بيا
ابتدايي منطقي بدلونونه ترسره کوو.

\begin{thm}[درې ګونتوب]\ollabel{ordtrichotomy}
$\alpha \in \beta \lor \alpha = \beta \lor \beta \in \alpha$، د هر دوو
ترتيبي عددونو $\alpha$ او $\beta$ دپاره.
\end{thm}

\begin{proof}
ثبوت په دوه ګونې استقرا ولاړ دے، يعنې
\olref{ordinductionschema} دوه ځله کاروو. وايو چې $x$ له $y$ سره
\emph{د پرتله کېدو وړ} دے، که او يوازې که $x \in y \lor x = y \lor y \in x$ وي.

د استقرا دپاره فرض کړئ چې په~$\alpha$ کښې هر ترتيبي عدد له
\emph{هر} ترتيبي عدد سره د پرتله کېدو وړ دے. د دويمې استقرا دپاره فرض کړئ
چې $\alpha$ په~$\beta$ کښې له هر ترتيبي عدد سره د پرتله کېدو وړ دے.
و به ښيو چې $\alpha$ له~$\beta$ سره د پرتله کېدو وړ دے. په~$\beta$ باندې
له استقرا څخه به نتيجه واخلو چې $\alpha$ له هر ترتيبي عدد سره د پرتله کېدو
وړ دے؛ او بيا په~$\alpha$ باندې له استقرا څخه چې \emph{هر} ترتيبي عدد له
\emph{هر} ترتيبي عدد سره د پرتله کېدو وړ دے، لکه چې پکار دي.
بس دا فرض کول کافي دي چې $\alpha \notin \beta$ او $\beta \notin
\alpha$ دي، او وښيو چې $\alpha = \beta$ دے.

د دې د ښودلو دپاره چې $\alpha \subseteq \beta$ دے، $\gamma \in \alpha$
ثابت وټاکئ؛ د \olref{ordmemberord} له مخې دا ترتيبي عدد دے. نو د لومړۍ
استقرا د فرض له مخې $\gamma$ له $\beta$ سره د پرتله کېدو وړ دے.
خو که يا $\gamma
= \beta$ يا $\beta \in \gamma$ وي، نو $\beta \in \alpha$ دے
(که اړتيا وي، د $\alpha$ له تعديت څخه کار اخلو)، چې زموږ د فرض خلاف دے؛
نو $\gamma \in \beta$ دے. په تعميمولو سره $\alpha \subseteq
\beta$ دے.

په عين ورته استدلال، د دويمې استقرا د فرض په کارولو، ښيو چې
$\beta \subseteq \alpha$ دے. نو $\alpha = \beta$ دے.
\end{proof}\noindent له دې امله، کله ناکله به د $\alpha <\beta$
نښه د $\alpha \in \beta$ پر ځای وليکو، ځکه چې $\in$ د ترتيب د اړيکې
په شان کار کوي. له اشناوالي پرته کومه ژوره وجه نۀ لري؛ او هم دا چې
$\alpha \leq \beta$ ليکل تر $\alpha \in
\beta \lor \alpha = \beta$ اسانه دي.\footnote{موږ $\alpha
\mathrel{\underline{\in}} \beta$ ليکلے شو؛ خو دا به بالکل غير معياري وي.}

زموږ د وروستيو نتيجو دوه لنډې پايلې دا دي؛ لومړۍ پايله نوې نښې کاروي:

\begin{cor}\ollabel{ordordered}
که $\exists \alpha\phi(\alpha)$ وي، نو $\exists \alpha(\phi(\alpha)
\land \forall \beta(\phi(\beta) \lif \alpha \leq \beta))$ دے.
همدارنګه، د هر درې ترتيبي عددونو $\alpha, \beta, \gamma$ دپاره دواړه
$\alpha
\notin \alpha$ او $\alpha \in \beta \in \gamma \lif \alpha \in
\gamma$ دي.
\end{cor}

\begin{proof}
د \olref[wo]{wo:strictorder} په شان.
\end{proof}

\begin{prob}
د \olref[sth][ordinals][basic]{ordtrichotomy} په ثبوت کښې «په عين ورته
استدلال» بشپړ کړئ.
\end{prob}

\begin{cor}\ollabel{corordtransitiveord}
$A$ ترتيبي عدد دے، که او يوازې که $A$ د ترتيبي عددونو متعدي سټ وي.
\end{cor}

\begin{proof}
\emph{له چپ نه ښي خوا ته.} د \olref{ordmemberord} له مخې.
\emph{له ښي نه چپ خوا ته.} که $A$ د ترتيبي عددونو متعدي سټ وي، نو
$\in$ د $A$ ښه ترتيب د \olref{ordinductionschema} او
\olref{ordtrichotomy} له مخې ورکوي.
\end{proof}

اوس، موږ د
\olref{ordinductionschema} او \olref{ordtrichotomy} معنا داسې بيان کړه چې
$\in$ د ترتيبي عددونو ښه ترتيب ورکوي. خو د لاندې نتيجې له امله بايد
د دې ډول ادعا په اړه \emph{ډېر احتياط} وکړو:

\begin{thm}[د بورالي--فورتي تناقض]\ollabel{buraliforti}
د ټولو ترتيبي عددونو سټ نشته.
\end{thm}

\begin{proof}
د تناقض د راوستلو دپاره فرض کړئ $O$ د ټولو ترتيبي عددونو سټ دے.
که $\alpha \in
\beta \in O$ وي، نو د \olref{ordmemberord} له مخې $\alpha$ ترتيبي عدد دے،
نو $\alpha \in O$ دے. په دې توګه $O$ متعدي دے، او له
\olref{corordtransitiveord} څخه $O$ ترتيبي عدد دے. نو $O \in O$ دے، چې له
\olref{ordordered} سره په تناقض کښې دے.
\end{proof}
\noindent
دا نتيجه د \citeauthor{Burali-Forti1897} په نوم نومول شوې ده. خو کانتور په
۱۸۹۹ کښې، ډېډېکنډ ته په يوه ليک کښې، په لومړي ځل په روښانه توګه وليدل
چې د ټولو ترتيبي عددونو د يوه سټ فرض کول \emph{تناقض} راولاړوي.
لکه څنګه چې فان هاينورټ وضاحت کوي:
\begin{quote}
  بورالي--فورتي پخپله دا تناقض داسې ګاڼه چې د \emph{برهان بالخلف}
  له لارې دا نتيجه ثابتوي چې د ترتيبي عددونو طبيعي ترتيب يوازې
  يو جزوي ترتيب دے.
  \citep[p.~105]{Heijenoort1967}
\end{quote}
د بورالي--فورتي تېروتنه يوې خوا ته پرېږدو؛ پورته بحث داسې لنډولے شو:
ترتيبي عددونه هغه سټونه دي چې هر يو په جلا توګه د غړيتوب له مخې ښه مرتب
دے، او ټول په ګډه د غړيتوب له مخې ښه مرتب دي، خو په ګډه يو سټ نۀ جوړوي.

د دې بحث په پای کښې، د ترتيبي عددونو ځينې نور بنسټيز خاصيتونه وړاندې
کوو چې له \olref{ordinductionschema} او
\olref{ordtrichotomy} څخه نتيجه اخلو.

\begin{prop}
د ترتيبي عددونو هره په کلکه نزولي لړۍ متناهي ده.
\end{prop}

\begin{proof}
د ترتيبي عددونو هره نامتناهي په کلکه نزولي لړۍ $\alpha_0 > \alpha_1 > \alpha_2 > \ldots$
د $<$ له مخې مينيمال غړے نۀ لري، چې له \olref{ordinductionschema} سره
په تناقض کښې ده.
\end{proof}

\begin{prop}\ollabel{ordinalsaresubsets}
$\alpha \subseteq \beta \lor \beta \subseteq \alpha$، د هر دوو ترتيبي
عددونو $\alpha, \beta$ دپاره.
\end{prop}

\begin{proof}
که $\alpha \in \beta$ وي، نو $\alpha \subseteq \beta$ دے، ځکه چې $\beta$
متعدي دے. همدارنګه، که $\beta \in \alpha$ وي، نو $\beta \subseteq
\alpha$ دے. او که $\alpha = \beta$ وي، نو $\alpha \subseteq \beta$ او
$\beta \subseteq \alpha$ دي. نو د \olref{ordtrichotomy} له مخې مقصد ترلاسه شو.
\end{proof}

\begin{prop}\ollabel{ordisoidentity}
$\alpha = \beta$، که او يوازې که $\ordeq{\alpha}{\beta}$ وي، د هر دوو
ترتيبي عددونو $\alpha, \beta$ دپاره.
\end{prop}

\begin{proof}
ترتيبي عددونه ښه ترتيبونه دي؛ نو دا له درې ګونتوب
(\olref[basic]{ordtrichotomy}) او
\olref[iso]{wellordnotinitial} څخه سمدستي نتيجه اخلو.
\end{proof}

\begin{prob}\ollabel{probunionordinalsordinal}
ثابت کړئ چې، که د $X$ هر غړے ترتيبي عدد وي، نو $\bigcup X$ ترتيبي عدد دے.
\end{prob}

\end{document}
